\documentclass[10pt,a4paper]{amsart}
\usepackage{amsmath,amssymb,amsthm,graphicx}
% For commutative diagrams we use Xy-pic
\usepackage[all]{xy}

% We use 2cell for 2-commutative diagrams.
\xyoption{2cell}
% Theorem environments.
%
\theoremstyle{plain}
\newtheorem{theorem}[subsection]{Theorem}
\newtheorem{proposition}[subsection]{Proposition}
\newtheorem{lemma}[subsection]{Lemma}

\theoremstyle{definition}
\newtheorem{definition}[subsection]{Definition}
\newtheorem{example}[subsection]{Example}
\newtheorem{exercise}[subsection]{Exercise}
\newtheorem{situation}[subsection]{Situation}

\theoremstyle{remark}
\newtheorem{remark}[subsection]{Remark}
\newtheorem{remarks}[subsection]{Remarks}

\numberwithin{equation}{subsection}

% Macros
%
\def\lim{\mathop{\mathrm{lim}}\nolimits}
\def\colim{\mathop{\mathrm{colim}}\nolimits}
\def\Spec{\mathop{\mathrm{Spec}}}
\def\Hom{\mathop{\mathrm{Hom}}\nolimits}
\def\Ext{\mathop{\mathrm{Ext}}\nolimits}
\def\SheafHom{\mathop{\mathcal{H}\!\mathit{om}}\nolimits}
\def\SheafExt{\mathop{\mathcal{E}\!\mathit{xt}}\nolimits}
\def\Sch{\mathit{Sch}}
\def\Mor{\mathop{\mathrm{Mor}}\nolimits}
\def\Ob{\mathop{\mathrm{Ob}}\nolimits}
\def\Sh{\mathop{\mathit{Sh}}\nolimits}
\def\NL{\mathop{N\!L}\nolimits}
\def\CH{\mathop{\mathrm{CH}}\nolimits}
\def\proetale{{pro\text{-}\acute{e}tale}}
\def\etale{{\acute{e}tale}}
\def\QCoh{\mathit{QCoh}}
\def\Ker{\mathop{\mathrm{Ker}}}
\def\Im{\mathop{\mathrm{Im}}}
\def\Coker{\mathop{\mathrm{Coker}}}
\def\Coim{\mathop{\mathrm{Coim}}}

% Boxtimes
%
\DeclareMathSymbol{\boxtimes}{\mathbin}{AMSa}{"02}

%
% Macros for moduli stacks/spaces
%
\def\QCohstack{\mathcal{QC}\!\mathit{oh}}
\def\Cohstack{\mathcal{C}\!\mathit{oh}}
\def\Spacesstack{\mathcal{S}\!\mathit{paces}}
\def\Quotfunctor{\mathrm{Quot}}
\def\Hilbfunctor{\mathrm{Hilb}}
\def\Curvesstack{\mathcal{C}\!\mathit{urves}}
\def\Polarizedstack{\mathcal{P}\!\mathit{olarized}}
\def\Complexesstack{\mathcal{C}\!\mathit{omplexes}}
% \Pic is the operator that assigns to X its picard group, usage \Pic(X)
% \Picardstack_{X/B} denotes the Picard stack of X over B
% \Picardfunctor_{X/B} denotes the Picard functor of X over B
\def\Pic{\mathop{\mathrm{Pic}}\nolimits}
\def\Picardstack{\mathcal{P}\!\mathit{ic}}
\def\Picardfunctor{\mathrm{Pic}}
\def\Deformationcategory{\mathcal{D}\!\mathit{ef}}
\newenvironment{reference}{}{}
\newenvironment{slogan}{}{}
\newenvironment{history}{}{}
\usepackage{cleveref}
\crefname{originallemma}{original lemma}{original lemmas}
\Crefname{originallemma}{Original lemma}{Original lemmas}
\makeatletter
\newlabel{lemma-divisor-special-fiber}{{0C5Z}{source}{}{sourcepointer-1}{}}
\newlabel{lemma-divisor-special-fiber@cref}{{[originallemma][0][]0C5Z}{[1][0][]source}{}{}{}}
\makeatother
\title{The divisor obtained by dividing the special fibre by the gcd of its multiplicities is torsion of order dividing that gcd}
\author{The Stacks Project Authors}
\date{stacks-project:models.tex:local-snapshot-20261012; GNU FDL 1.2 or later}
\begin{document}
\maketitle
\noindent\textit{Editorial note.} The counted unit is the Stacks Project Lemma lemma-multiple-fibre-normal-bundle at source lines 4792-4799 together with its proof at source lines 4801-4810, a straight adjacent pairing in the chapter file. The chapter is the upstream file models.tex, which is the file name recorded in the source record; the pool directory carries the longer name stacks-models and is not used as the chapter name anywhere in this card.\par\medskip
\noindent\textit{Editorial note.} The retained statement opens with the words In Situation followed by a reference to situation-regular-model, so the situation environment that defines the setting is retained as a context span at source lines 4628-4637. It is context, not a counted unit, and the same context span is retained by two other cards of this chapter; no counted interval is shared.\par\medskip
\noindent\textit{Editorial note.} This chapter's candidate index records the statement body and the proof body without their closing markers, so the retained spans used here are the full environments, exactly as in the earlier cards of this chapter.\par\medskip
\noindent\textit{Editorial note.} The retained proof quotes one lemma of the same chapter, once through the retained situation and once in its own first line; it is declared as a single non-load-bearing navigation pointer with its stable Stacks tag looked up in the shared tag index. No citation occurs in the retained spans, so no bibliography entry is supplied.\par\medskip
\noindent\textit{Editorial note.} The retained preamble is the chapter preamble's Xy-pic setup at lines 14-18, the theorem declarations and mathematical macros at lines 154-217, and the three no-op wrapper environments.\par\medskip
\noindent\textit{Allowed standard prerequisites (editorial).}\begin{itemize}
\item The regular-model situation, the special fibre as a divisor, and the conormal sheaf of a divisor. In the regular-model situation the special fibre is written as the sum of its irreducible components with multiplicities and is an effective Cartier divisor; for an effective Cartier divisor D on a scheme X the conormal sheaf of D in X is the pullback of the invertible sheaf attached to the negative of D. The author's own situation and his lemma identifying the special fibre as a Cartier divisor are the allowed inputs.. The counted lemma produces the torsion statement that the chapter uses to control the Picard group and the normal bundle in the analysis of multiple fibres.
\end{itemize}\medskip
\ifdefined\hypertarget\hypertarget{sourcepointer-1}{}\fi
\noindent\textit{Source reference.} Original lemma 0C5Z: focused-seeds/AGP/sources/stacks-snapshot-20261011/tags.txt, source lines 15279--15279. Navigation pointer inside the retained situation and invoked again by the retained proof for the identification of the special fibre as an effective Cartier divisor; that identification is an allowed input. The number is the stable Stacks tag of the target, looked up in the shared tag index.\par\medskip
\begin{situation}
\label{situation-regular-model}
Let $R$ be a discrete valuation ring with fraction field $K$,
residue field $k$, and uniformizer $\pi$.
Let $C$ be a smooth projective curve over $K$ with $H^0(C, \mathcal{O}_C) = K$.
Let $X$ be a regular proper model of $C$.
Let $C_1, \ldots, C_n$ be the irreducible components of the special
fibre $X_k$. Write $X_k = \sum m_i C_i$ as in
Lemma \ref{lemma-divisor-special-fiber}.
\end{situation}

\begin{lemma}
\label{lemma-multiple-fibre-normal-bundle}
In Situation \ref{situation-regular-model} set $d = \gcd(m_1, \ldots, m_n)$
and let $D = \sum (m_i/d)C_i$ as an effective Cartier divisor.
Then $\mathcal{O}_X(D)$ has order dividing $d$ in $\Pic(X)$
and $\mathcal{C}_{D/X}$ an invertible $\mathcal{O}_D$-module
of order dividing $d$ in $\Pic(D)$.
\end{lemma}

\begin{proof}
We have
$$
\mathcal{O}_X(D)^{\otimes d} = \mathcal{O}_X(dD) =
\mathcal{O}_X(X_k) = \mathcal{O}_X
$$
by Lemma \ref{lemma-divisor-special-fiber}.
We conclude as $\mathcal{C}_{D/X}$ is the pullback of
$\mathcal{O}_X(-D)$.
\end{proof}

\end{document}
